\section{Theoretical background}
Contrary to regular filtering, the designer targets a specific frequency or bandwidth, 
placing zeroes and poles on the unit-circle to \textit{approximate}, a desired result, or 
atleast get somewhat close to something acceptable.\\

\subsection{LMS}
In adaptive filtering, we're more concerned about the system as an input and how we implement
an algorithm to force the error to be zero. In this lab journal the overall system we'll be examining
will take the form of:

\begin{figure}[htpb]
    \centering
    \begin{tikzpicture}[
        >=Latex,
        thick,
        block/.style={
            draw,
            rectangle,
            minimum width=2.5cm,
            minimum height=1.2cm,
            align=center
        },
        sum/.style={
            draw,
            circle,
            minimum size=0.8cm,
            inner sep=0pt
        }
    ]
    \node[block] (system) at (4,2) {system};
    \node[block] (fir) at (4,0) {FIR filter};
    \node[block] (adaptive) at (4,-2)
        {adaptive\\algorithm};
    
  % Error summation junction
    \node[sum] (sum) at (7,2) {$+$};

    % Main input line to system
    \draw[->] (0,2) -- node[above] {$x[k]$} (system.west);

    % Branch downward from main input line
    \draw (1.5,2) -- (1.5,-2);

    % Branches to FIR and adaptive algorithm
    \draw[->] (1.5,0) -- (fir.west);
    \draw[->] (1.5,-2) -- (adaptive.west);
    \draw[->] (system.east) --
        node[above] {$d[k]$} (sum.west);

    % Branch combined input to FIR and adaptive algorithm
    \draw[->] (1.5,2) |- (fir.west);
    \draw[->] (1.5,0) |- (adaptive.west);

    % System output to error summation
    \draw[->] (system.east) --
        node[above] {$d[k]$} (sum.west);

    % FIR output to error summation
    \draw[->] (fir.east) -| (7,1.6);
    \node at (6.1,0.25) {$y[k]$};
    \node at (7.35,1.35) {$-$};

    % Error output
    \draw[->] (sum.east) --
        (9,2) node[above left]
        {$e[k]=d[k]-y[k]$};

    % Error feedback to adaptive algorithm
    \draw[->] (9,2) |- (adaptive.east);

    % Adaptive algorithm updates FIR coefficients
    \draw[->] (adaptive.north) -- (fir.south);


    \end{tikzpicture}
    \caption{Adaptive FIR filter structure.}
    \label{fig:adaptive_filter}
\end{figure}

Where $x[k]$ will be some randomly generated noise, with a zero-mean, unit variance, passing through a system
summed together with a sinusoid $s[k]$, for testing purposes, to experimentally see how the system responds
to both noise and a signal combined.\\

In this example, the unidentified system will be $d[k]= 0.5x[k] + 0.25x[k-1] -0.3x[k-2]$\\

The simple LMS filter works by taking the output of the system, as the desired output $d[k]$ and then
subtracting $y[k]$, which is the output of the FIR filter, governed by a vector of weights $w[k]$. The
The LMS algorithm will then try to force the error $e[k] = d[k] - y[k] \approx 0$ by tuning the weights 
in descrete steps.

The vector $\mathbf{W}[k]$ hold the coefficients and the vector $\mathbf{X}[k]$ holds old samples :
\[
\mathbf{W}[k] = 
\begin{bmatrix}
    w_0[k]\\
    w_1[k]\\
    ...\\
    w_N-1[k]\\
\end{bmatrix},
\qquad
\mathbf{X}[k] = 
\begin{bmatrix}
    x[k]\\
    x[k-1]\\
    ...\\
    x[k-N+1]\\
\end{bmatrix}
\]

The complete algorithm for computing the coefficiens for the vector:
\[\mathbf{W}[k+1] = \mathbf{W}[k] + 2\mu e[k]\mathbf{x}[k]\]
The factor $\mu$ is factor that gives some control of how aggressize the filter should act, as it a scalar
that is multiplied onto the weights and the old samples. The idea of the algorithm is then, that each iteration
that it runs, the weights of the filter are multipled with the signal, which affects how $y[k+1]$ samples will behave,
this affects the error, and it repeats. Each iteration the weights gets updates aswell, so that they are
applied with a different value at the next iteration. 


\subsection{Variable Step-size LMS (VVS LMS)}
As briefly mentioned in the LMS section, the factor $\mu$ is a deciding factor
in filter performance for an LMS filter. It determines the speed of convergence and accuracy.
In the variable step-size algorithm each filter tab has an independent.
It follows the form:
\begin{equation}
% w_i[n+1] = w_i[n] + 2\mu_i[n]e[n]x[n-i],
w_o[k+1] = w_i[k] + 2\mu_i[k]g_i[k]
,\qquad
\textit{for} \qquad i = 0, 1, \dots, N-1
\label{VSS_LMS_W}
\end{equation}
  
The interesting part of this implementation is the gradient term, which
is an intermediate calculation for the update of the $\mu$ value.
\begin{equation}
g_i[k] = e[k]x[k-i] \qquad  \qquad 
\label{VSS_LMS_g}
\end{equation}
\textit{and}
\begin{equation}
u_i[k] = u_i[k-1] + \rho sign \left[g_i[i]\right]sign\left[lg_i[k-1]\right]
\label{VSS_LMS_mu_update}
\end{equation}
It is then up to the individual use case to determine a $\mu_{min}$ and $\mu_{max}$
This essentailly calculates a gradient for direction, positive or negative, that the $\mu$ should move, granting
the ability to increase or decrease the value. Another tuneable factor is $\rho$ which essentially dictates how much 
the step-size itself changes, as it sits as an index independt factor, just as $\mu$ did in the regular LMS.
